\cotefile{51}
\scanpage{2}
``Bonne réduction''

(\unsure{Shimura|shimura}, Koizumi)
% archive sheet number: 1

\scanpage{3}
\noindent\hfill\fbox{\unsure{Shimura} - Koizumi}

\uline{Proposition} 1 Soient $V$ un anneau de valuation
discrète, \inserted{de corps $K$} $X$ un \unsure{préschém.|préschéma} sur $V$, $F$ un faisceau
quasi-cohérent sur $X$, \inserted{$X'$ l'\unsure{ouvert} $\unsure{X_K|X}$ de $X$, $F'=F_K=F\vert{}X'$}
\deleted{\unsure{Pour} \unsure{que|qu.} \unsure{quotient} \unsure{Pour} que $F$ soit plat sur
$V$, il f. et s. que le plus petit \unsure{de schém.|sous-schéma}
\unsure{quo.|q.} coh. $G$ de $F$ soit plat sur $V$, il \unsure{f}}

\deleted{\unsure{À t}} (i) Pour tout faisceau quotient q. coh.
$G'=F'_{\deleted{\unsure{K}}}/J'$, soit $G$ le faisceau
quotient \deleted{de} $F/J^*$, où $J^*$ est le plus
\deleted{\unsure{grand|plus} petit \unsure{faisceau}}\inserted{grand ss-Module} de $F$ induisant \deleted{$J'$}
\deleted{sur} sur $X'=X_K$ un \deleted{\unsure{sous}} Module $\subset J'$, --
donc $G$ est le plus petit \inserted{Module} quotient de $F$
induisant sur \unsure{$\vert$} $X'_{\deleted{\unsure{K}}}$ un \deleted{\unsure{sous}} faisceau
quotient de $F'$ \deleted{plus grand} majorant $G'$.
Alors $G$ est quasi-cohérent et $V$-plat, et
\[
G|X'=G'
\]

(ii) \deleted{Soit $G=F/$ \unsure{Le} \unsure{un}} Soit $G=F/J$ un
faisceau q. coh de \deleted{\unsure{$G$|$F$}}$F$ qui est plat sur $V$,
et soit $G'=G|X'$ le \unsure{qu.|quotient}\deleted{\unsure{t}} correspondant
de $F'$. Alors $G$ est le plus petit\deleted{\unsure{e}}
\deleted{\unsure{M.}}\unsure{ent} quotient de \deleted{\unsure{$G$}}$F$ \deleted{\unsure{majorant}} qui
induise $G'$ sur $X'$.

\uline{Corollaire 1} Correspondance biunivoque entre
l'ensemble des faisceaux q. coh. \deleted{\unsure{plats}} de
$F$ qui sont $V$-plats, et des faisceaux \inserted{quotients de $F'$} q. coh.
% archive sheet number: 2

\scanpage{4}
\uline{Dém} On sait déjà (\unsure{réf|bof}) que $J$ est
q. coh. et $J\vert{}X'=J'$, donc $G$ est q. \unsure{coh|ch}
et $G|X'=G'$. Prouvons que $G$ est $V$-plat,
i.e. que l'uniformisante $\pi$ de $V$ est
non diviseur de $0$ dans $G$. \deleted{\unsure{Sinon,}}
Soit en effet \deleted{\unsure{$J$} \unsure{$K$}} $\unsure{\mathcal N|N}$ \deleted{\unsure{l'élément}} noyau
de $\pi$ dans $G$, alors $G/\mathcal N$ est un
\deleted{faisceau}\inserted{Module} quotient \deleted{\unsure{qu}} de $G$ qui induit
encore $G'$ sur $X'$, donc d'après la
condition minimale de $\deleted{\unsure{\mathcal N}}G$, \deleted{\illegible{1}} on a $\mathcal N=0$,
ce qui prouve (i).

Pour (ii), on note que si
$G_1$ est le plus petit \unsure{module|Module} quotient
de $G$ induisant $G'$ sur $X'$, \deleted{\unsure{Alors}}
de sorte que $G_1$ est \inserted{un quotient de $F$} majoré par $G$,
\margin{$G\to G_1$}alors le noyau de \inserted{\unsure{l'épim|libérer}} $G\to G_1$ est \inserted{\unsure{encore|aussi|inclus}} plat sur $V$,
et induit $0$ sur $X'$, il suffit de
prouver qu'un tel Module q. coh. est
Nul. \unsure{Cela|Or} est immédiat, puisque
\[
N\otimes_V K=0.\qquad \text{Cqfd.}
\]
% archive sheet number: 3

\scanpage{5}
\uline{Corollaire 2} Supposons qu'on ait deux
$F_1$ \unsure{resp$^t$} $F_2$ \inserted{q. coh.} sur $X_1$ \inserted{\unsure{resp|resp.} $X_2$}, et qu'on considère les
\deleted{faisceaux}\inserted{Modules} quotients de $F_1$, $F_2$, $F_1\otimes_V F_2$. Alors
la correspondance du Cor. 1 est compatible
avec la formation du \deleted{\unsure{pro} \unsure{f}} produit tensoriel
de quotients des \deleted{\unsure{$F_i$|$F_1$}, \unsure{faisceaux} \unsure{quotients}}\inserted{$F_1$ \deleted{\unsure{et}} et $F_2$} ($F'_1$ et $F'_2$)
identifié à un quotient de $F_1\otimes F_2$ ($F'_1\otimes F'_2\unsure{\neq|{+}}$)

En effet, le produit tensoriel \uline{sur $V$}
de faisceaux qu. coh. $V$-plats est $V$-plat,
et la \deleted{\unsure{pass}} formation du produit tensoriel
commute à la tensorisation par $K$.

\uline{\deleted{\unsure{P}}Corollaire 3} \deleted{Supposons que} Les \unsure{souspréschémas|sous-préschémas}
fermés $Y$ de $X$ plats sur $\unsure{N|V}$ sont en
correspondance biunivoque avec les
sous-préschémas fermés $Y'$ de $X'=X\otimes_V K$,
à $Y$ correspond $Y'=Y\otimes_V K=Y\cap X'$, et
à $Y'$ \unsure{correspond|correspondent} le sous-préschéma
adhérence $Y=\overline{Y'}$. Cette correspondance
est compatible avec les produits
cartésiens sur $V$.
% archive sheet number: 4

\scanpage{6}
\uline{Corollaire 4} Soit $G$ un préschéma en
groupe sur $V$, $H'$ un sous-préschéma en
groupe \inserted{fermé} de $G_K$, $H$ le sous-préschéma
adhérence de $H'$. Alors $H'$ est un
sous-préschéma en groupe de $G$, plat
sur $V$, \deleted{dominant} induisant $H'$, et c'est
le seul.

Il est \unsure{trivial|évident} que $H$ est stable par
\unsure{symmétrie|symétrie}, reste à prouver que
$\unsure{p|\rho}:G\times G\mathrel{-}G$ induit $H\times H\mathrel{-}H$, i.e.
que le sous-schéma \inserted{fermé} $P$ de $G\times G$
image inverse de $H$ par $p$ \deleted{de} majore
$H\times H$. Or ce sous-schéma induit
sur $X'$ un sous-schéma qui majore
$H'\times H'$, \deleted{donc} il majore $\overline{H'\times H'}$ qui est
$H\times H$ d'après le cor. 3. Cqfd.
% archive sheet number: 5

\scanpage{7}
\uline{Prop. 2}

Soit $G$ un schéma en groupes sur $S$ loc noeth,
on suppose que $\text{\uuline{Hilb}}_{\inserted{G/S}}\deleted{(G/S)}$ existe. On s'intéresse
% slanted
\margin{$G$ \unsure{qu.|q.} projectif sur $S$ ???}\unsure{à ens|à l'ens|l'ens|à l'ens.} $F(S)\subset\mathrm{Hilb}(G/S)$ formé des
sous-schémas fermés de $G$, propres et plats sur $S$,
qui sont des \uline{schémas \unsure{de|en} ss-groupes}. Pour tout
$S'$, sur $S$, on définit de $=$ $F(S')\subset\mathrm{Hilb}(G'/S')$.
d'où un ``sous foncteur'' de $\text{\uuline{Hilb}}_{G/S}$.

\uline{Je dis qu'il est représentable} et que
\uline{le préschéma} $M$ \uline{sur} $S$ \uline{est tel que le
morphisme induit} $M\mathrel{-}\text{\uuline{Hilb}}_{G/S}$ \uline{est
une immersion fermée}. \uline{Si} $G$ est
\uline{propre sur} $S$ \uline{et} $\mathrm{Hilb}_{G/S}$ \uline{est réunion} disjointe
\uline{d'ouverts} \deleted{\unsure{qui} \unsure{pr}}\inserted{$H_i$} \uline{de type fini \deleted{\unsure{$H_i$}} sur} $S$,
\deleted{\unsure{alors} les $H_i\cap M=M_i$} \inserted{donc} $M$ \unsure{--|est}
\uline{réunion disjointe des ouverts} $M_i=M\cap H_i$,
\uline{alors les} $M_i$ \uline{sont propres sur} $S$,

N.B. Si on s'intéresse à la \unsure{classification|caractérisation} des \inserted{ss} schémas \deleted{\unsure{s}}\inserted{en \unsure{grpes|groupes}} $H$
de $G$ qui sont \inserted{de \uline{plus} (déjà propres !)} finis (\unsure{resp$^t$|resp.} séparables, resp$^t$ abéliens),
on trouve des ss-schémas \uline{ouverts} de $M$, les premiers
\unsure{étant} également \uline{fermés}.
% archive sheet number: 6

\scanpage{8}
Les deux premières assertions résultent
du fait suivant : \uline{Soit} $H$ \uline{un sous-schéma
fermé de} $G$, \inserted{\uline{propre et}} \uline{plat sur} $S$. \uline{Alors il existe un}
\deleted{\uline{plus} \uline{petit}} \uline{sous-schéma fermé} $\deleted{S'}T$ \uline{de} $S$
tel que pour tt morphisme $\deleted{\unsure{T}}S'\mathrel{-}S$,
l'image inverse $H'$ est un ss-groupe de $G'$
\unsure{sss|ssi} $S'\mathrel{-}S$ est majoré par $T$. Or,
il faut exprimer, \deleted{si \unsure{$H$ et|$H$ est}} $1^\circ$) Les images
inverses de $H$ et $\check H$ sont les ${=}_{s}$
$2^\circ$) \deleted{\unsure{Les|l'} images}\inserted{L'image} inverse de $H\times H$ est
majorée par celle de $P$, où
$P=p^{-1}(H)$, $p:G\times G\mathrel{-}G$ est la
multiplication. Les deux \deleted{\unsure{que}} conditions
en question, il est connu, s'expriment
par la \deleted{\unsure{s}} domination de $S'\mathrel{-}S$ par
des ss-schémas fermés $T_1$, $T_2$ de $S$,
on prendra donc $T=T_1\unsure{\wedge|\cap} T_2$.

Pour la dernière assertion, le critère
habituel \unsure{des|de} propreté \unsure{nous} ramène au
Cor. 4 \unsure{,} de prop. 1. O.K.
% archive sheet number: 7

\scanpage{9}
\deleted{\uline{Corollaire}}

\uline{Proposition 3} Soit $A$ un schéma abélien sur $V$,
et $u':A':A_K\longrightarrow B'$ une isogénie, où $B'$ est
un schéma abélien sur $K$. Il existe alors
un schéma abélien $B$ sur $V$, \deleted{\unsure{se réduisant à} \unsure{un}}
tel que $B\otimes_V K\simeq B'$, et un \deleted{\unsure{en plus|et cette} isogénie}
ho\inserted{\unsure{mo}}morphisme \deleted{$A\mathrel{-}$} de schémas abéliens $u:A\mathrel{-}B$
induisant $u':A'\mathrel{-}B'$.

N.B. On ne se sert que du th. de Weil,
disant qu'un $V$-morphisme \deleted{\unsure{$A\mathrel{-}B$|$A_K\to B_K$} \unsure{entre}}
$A'\mathrel{-}B'$ \unsure{provient} d'un morphisme $A\mathrel{-}B$.

\uline{Dém}. Soit $H'$ le noyau de $u'$,
c'est un sous-groupe fini de $A'$, donc
provient d'un sous-groupe $V$-plat $H=\overline{H'}$ de
$A$, nécessairement \uline{fini} sur $V$, \unsure{étant} propre \inserted{et plat} sur
% slanted
\margin{\unsure{référence|retenons}.}$V$ à fibre générale finie. Comme
$A$ est \uline{projectif} sur $V$ (\unsure{réf}), on
peut construire le \inserted{schéma} quotient $A/\unsure{\Gamma|H}$,
qui est un schéma en groupes
\unsure{car|sur} $A$ est plat sur $V$, \uline{plat} sur $V$,
($=$ raisons), propre sur $V$, et tel que
% archive sheet number: 8
\scanpage{10}
le noyau de $A\longrightarrow B$ soit précisément $\unsure{\Gamma|H}$,
\deleted{\unsure{Pour|Prouvons} Montrer que $B$ est un schéma abélien}
\deleted{\unsure{Un}} Le passage au quotient en question
commute avec le changement de base, donc
$1^\circ$) $B\otimes_V K\simeq\deleted{\unsure{B}}A'/\Gamma'\simeq B'$ \qquad $2^\circ$) $B\otimes_V k$ est \deleted{\unsure{un}}
isom à \deleted{\unsure{un}} $A\otimes k/\Gamma'\otimes k$, donc est un
schéma abélien. Donc $B$ est un
schéma abélien. On a fini.

\uline{Corollaire} 1 \deleted{\unsure{Toute|Soit}} Si $A'$ \deleted{\unsure{est un}}\inserted{et $B'$ sont des}
schémas abéliens \inserted{isogènes} sur $K$, alors $A'$ ``se
réduit bien pour $V$'' sss $B'$ se réduit
bien pour $V$. S'il en est ainsi, les variétés
\deleted{\unsure{rédu}} abéliennes \unsure{réduites|induites} sont isogènes.

\deleted{\uline{Corollaire \unsure{2|3}}} Conséquence : On peut dire d'une
classe de variétés abéliennes mod isogénie
(p.ex. simple) si elle se réduit bien ou
non, et \deleted{\unsure{dire} \unsure{la}} parler de la
\deleted{variété \unsure{rédu}} classe réduite. Bien
entendu, une somme de classes
de \deleted{variétés}\inserted{schémas} abéliens sur $K$ qui se
réduisent bien, se réduit bien. En
particulier, si les composantes simples de
% archive sheet number: 9
\scanpage{11}
$A'$ se réduisent bien, $A'$ se réduit bien. On
sait prouver la réciproque, qui \unsure{équivaut|s'exprime}
au

\uline{Théorème} (\unsure{Shimura} - Koizumi) Si $A'$ sur $K$
se réduit bien, tt schéma abélien
quotient aussi.
% slanted
\margin{On peut le prouver \uline{directement} comme la proposition 3}

On \deleted{peut \unsure{prouver}} suppose $A'=A\otimes_V K$,
et le quotient en question \uline{facteur} direct.
On a donc des endomorphismes
$p'\unsure{,|;}q'$ de $A'$ tels que
\[
p'\unsure{+}q'=\mathrm{id},\qquad p'{}^2=p',\qquad q'{}^2=p'
\]
Ces endomorphismes correspondent à
des endomorphismes $p\unsure{,|;}q$ de $A$, \unsure{satisfaisant|satisfont}
les $\unsure{\mathrm m^{es}}$ conditions.
\[
p+q=\mathrm{id},\qquad p^2=p,\quad q^2=q.
\]
Donc on aura un isom
\[
A=\mathrm{Ker}\,p\times\mathrm{Ker}\,q=A_1\times A_2
\]
% margin note written sideways
\margin{\deleted{\unsure{Condition|Conclusion|Corollaire}}}(\unsure{c'est} ``fonctoriel'' !). Donc $A_1$ et
$A_2$ sont propres sur $V$, d'ailleurs la
somme de leurs fibres en $k$ est $A\otimes k$,
qui est \deleted{\unsure{simple}}\inserted{abélien} sur $k$, il faut que $A_1\otimes k$ et $A_2\otimes k$
soient abéliens sur $k$, donc \unsure{aussi} $A_1$ et $A_2$. D'où \unsure{le|la}
% archive sheet number: 10
\scanpage{12}
conclusion.

\uline{\deleted{\unsure{N}}Remarque} Soit $A$ un schéma abélien sur $V$,
$B'$ un sous-schéma abélien de $A'=A_K$,
Il résulte de ce qui précède que $B'$
est de la forme $B_K$, où $B$ est un
schéma abélien sur $V$, et où l'immersion
$B'\mathrel{-}A'$ est \unsure{induit|induite} par un morphisme
$B\longrightarrow A$ (ce dernier fait résulte
aussi du th. de Weil). \deleted{Mais} Soit
$H$ le sous-schéma fermé de $A$ image de $B$,
on voit \deleted{\unsure{facilement}}\inserted{tout de suite} que $H=\overline{B'}$ (car $H$
induit $B'$, \deleted{et} donc $H$ majore $\overline{B'}$,
or \deleted{\unsure{$B'$|$B$}} \unsure{l'|$1^\circ$}image inverse de $\overline{B'}$ est un
sous-schéma fermé de $B$ \unsure{donnant|dominant} $B\otimes K$, donc
identique à $B$, donc $\overline{B'}$ \deleted{est} majore $H$ :).
C'est donc un sous-schéma en groupes
fermé de $A$, \inserted{plat sur $V$.} (prop. 1, cor 4). C'est
un schéma abélien sss $B\longrightarrow B'$ est
un isom. (grâce au th de Weil p.ex),
i.e. $B\longrightarrow A$ est un monomorphisme,
i.e. son \uline{noyau} est \deleted{\illegible{1}} trivial. --
D'après \unsure{Shimura} - Koizumi, il n'en est
% archive sheet number: 11

